\section{Proof of the example}\label{sec:appendix}
\begin{proof}[Proof of Theorem~\ref{thm:geometric}]
Write $S_n=\sum_{k=0}^{n}r^k$. Multiplying by $r$ shifts each term by one index:
\[
rS_n=\sum_{k=1}^{n+1}r^k.
\]
Subtracting the second sum from the first cancels the terms with indices
$1,\ldots,n$, leaving $(1-r)S_n=1-r^{n+1}$. Since $r\ne1$, division by $1-r$
gives the claimed formula.
\end{proof}
